\section{Algorithm}\label{sec:algorithm}

\subsection{Augmented Lagrangian framework}\label{subsec:al_framework}

We formally state the first assumptions about problem~\eqref{pb:cnls}.
\begin{assumption}\label{assumption:functions_C2}
	Residuals $r$ and constraints $c$ are twice continuously differentiable on $\RR^n$.
\end{assumption}
\begin{assumption}\label{assumption:full_rank_A}
	The equality constraint matrix $A$ is full row rank.
\end{assumption}
\begin{assumption}\label{assumption:feasible_linear_cons}
	The feasible region for the linear constraints $\Omega$ is non-empty.
\end{assumption}

Our method is based on the Augmented Lagrangian (AL) function defined as
\begin{equation}\label{eq:al}
	\Phi(x,\lambda,\mu) := \dfrac{1}{2}\|r(x)\|^2 + \inner{\lambda}{c(x)} + \dfrac{\mu}{2} \|c(x)\|^2,
\end{equation}
where $\mu > 0$ is a penalty parameter.

We maintain the linear constraints in the formulation of the problem and only penalize the violation of the nonlinear constraints.
One has the following expression for the gradient: 
\begin{equation*}
	\nabla_x \Phi(x,\lambda,\mu) = J(x)^\top r(x) + C(x)^\top\bar{\lambda}(x,\lambda,\mu),
\end{equation*}
where
\begin{equation}\label{eq:1st_order_mult}
	\bar{\lambda}(x,\lambda,\mu):=\lambda + \mu c(x),
\end{equation} 
are the first-order estimates of the Lagrange multipliers. The AL and Lagrangian gradients are related by the identity
\begin{equation}\label{eq:identity_al_lag_grad}
	\nabla_x \Phi(x,\lambda,\mu) = \nabla_x \calL(x,\bar{\lambda}(x,\lambda,\mu)).
\end{equation}
The Hessian is given by
\begin{equation*}
	\nabla^2_{xx} \Phi(x,\lambda,\mu) = J(x)^\top J(x) + \mu C(x)^\top C(x) +  S(x) 
\end{equation*}
with the second-order terms explicitly given by
\begin{equation}\label{eq:al_hessian_2nd_terms}
	S(x) =\sum_{i=1}^{n_r} r_i(x) \nabla^2r_i(x) + \sum_{i=1}^{n_c} \nabla^2 c_i(x) \bar{\lambda}_i(x,\lambda,\mu).
\end{equation}
For fixed $\lambda$ and $\mu$, reformulating problem~\eqref{pb:cnls} with function~\eqref{eq:al} gives the linearly constrained problem
\begin{equation}\label{pb:cnls_al_reformulation} 
	\begin{aligned}
		\min_{x} \quad& \Phi(x,\lambda,\mu)  \\
		\text{s.t.}  \quad & x \in \Omega.
	\end{aligned}	
\end{equation}
Moving the nonlinear constraints into the objective simplifies the problem because only linear constraints are left, which enables one to use iterative methods for linearly constrained optimization while still improving feasibility of the nonlinear constraints. Since $\Omega$ is a convex set, a first-order critical point of~\eqref{pb:cnls_al_reformulation} can also be characterized by the condition
\begin{equation*}
	\proj{\Omega}{x^*-\nabla_x \Phi(x^*,\lambda,\mu)} = x^*.
\end{equation*}
In our methods, and as in standard AL methods, each iterate $x_k$ is an approximate solution of a subproblem of the form~\eqref{pb:cnls_al_reformulation} satisfying
\begin{equation}\label{eq:algo_traulls_criticality_test}
	\|\proj{\Omega}{x_k-\nabla_x\Phi(x_k, \lambda_k, \mu_k)}-x_k\| \le \omega_k,
\end{equation}
for a fixed vector of multipliers $\lambda_k$, penalty parameter $\mu_k$ and tolerance $\omega_k$.
If the iterate $x_k$ improves the feasibility violation of the nonlinear constraints, a dual ascent step is taken on the Lagrange multipliers by setting
\begin{equation*}
	\lambda_{k+1} = \lambda_k + \mu_k c_k = \bar{\lambda}(x_k,\lambda_k,\mu_k).
\end{equation*}
If on the contrary, feasibility has not been improved, the minimization~\eqref{pb:cnls_al_reformulation} is restarted using an increased penalty parameter $\mu_{k+1} > \mu_k$ to enforce feasibility. The method is outlined in Algorithm~\ref{algo:traulls}. Following our notation for iteration dependent quantities, we respectively write $\nabla_x\Phi(x_k, \lambda_k, \mu_k)$ and $\bar{\lambda}(x_k,\lambda_k,\mu_k)$ in the compact form $\nabla_x \Phi_k$ and $\bar{\lambda}_k$.
\begin{algorithm}
	\caption{Augmented Lagrangian algorithm}\label{algo:traulls}
	\begin{algorithmic}[1]
		\Require Initial starting point $x^s_0 \in \Omega$, Lagrange multipliers estimate $\lambda_0$, penalty parameter $\mu_0$
		\State Penalty parameter increase factor $\tau > 1$, tolerance update constants $\omega,\eta, \kappa_\omega, \kappa_\eta, \beta_\omega, \beta_\eta$
		\State{Set initial tolerances $\omega_0\gets \omega \mu_0^{-\kappa_\omega}$ and $\eta_0\gets \eta \mu_0^{-\kappa_\eta}$}
		\For{k=0,1,2\ldots}
		\State\label{line:inner_iteration} Starting from $x_k^s$, approximately solve $\min_{x\in \Omega} \Phi(x,\lambda_k,\mu_k)$ to find $x_k \in \Omega$ satisfying~\eqref{eq:algo_traulls_criticality_test}.
		\If{$\|c(x_k)\| \le \eta_k$}
		\If{\(\|\proj{\Omega}{x_k-\nabla_x\Phi_k}-x_k\|  \le \omega_*\) and \(\left\Vert c(x_k)\right\Vert \le \eta_*\)}
		\State{\textbf{stop} and} \Return approximate solution $(x_k,\lambda_k)$
		\EndIf
		\State Update the Lagrange multipliers $\lambda_{k+1}\gets \bar{\lambda}_k$ 
		\State Set the next starting point $x_{k+1}^s \gets x_k $
		\State Leave the penalty parameter unchanged \(\mu_{k+1} \gets \mu_k\)
		\State Decrease the tolerances $\omega_{k+1}\gets \omega_k \mu_{k+1}^{-\beta_\omega}$ and $\eta_{k+1}\gets \eta_k \mu_{k+1}^{-\beta_\eta}$
		\Else{}
		\State Let the iterate unchanged \(\left(x_{k+1}^s,\lambda_{k+1}\right) \gets \left(x_k^s,\lambda_k\right)\)
		\State Increase the penalty parameter $\mu_{k+1} \gets \tau\mu_k$
		\State Update tolerances $\omega_{k+1}\gets \omega \mu_{k+1}^{-\kappa_\omega}$ and $\eta_{k+1}\gets \eta \mu_{k+1}^{-\kappa_\eta}$
		\EndIf
		\EndFor
	\end{algorithmic}
\end{algorithm}
One first notices that Algorithm~\ref{algo:traulls} has an outer-inner iteration structure, in the sense that each iteration requires to approximately solve an optimization problem (line~\ref{line:inner_iteration}), which also involves an iterative process. The techniques used for the latter are detailed in Section~\ref{sec:inner_iteration}.

Tolerances in Algorithm~\ref{algo:traulls} are updated in such a way that both sequences $\omega_k$ and $\eta_k$ tend to $0$ when $k \to \infty$. The update scheme that we outline follows the rules described in~\cite[][Chapter 14]{conn-etal:2000}. The parameters values used in our implementation are given in Section~\ref{subsec:implementation_details}.

We describe the convergence in terms of the norm of the projected gradient but this quantity is not easily computable in practice. Indeed, contrary to the case where there is only bounds or linear equality constraints, there is no direct formula for the projection on a polyhedral set of the form $\Omega$. For our implementation, we will use the norm of the reduced gradient, which we define now. 

For $x \in \Omega$, $I^+(x)$, resp. $I^-(x)$ denotes the indices of upper bounds, resp. lower bounds, satisfied as equalities (active) at $x$ and $I(x) := I^+(x) \cup I^-(x)$, i.e.:
\begin{equation*}
	i \in I(x) \implies x_i \in \{l_i,u_i\}.
\end{equation*}
We also introduce the notion of tangent space at a point $x$, written $T(x)$ and defined as the set of directions $d$ such that the same constraints are satisfied with equality at both $x$ and $x+d$. In our case, this corresponds to the linear equality constraints and the active bounds. Formally:
\begin{equation}\label{eq:tangent_space}
	T(x) := \left\{ d \in \RR^n \ | \ Ad = 0,\ d_i=0\ \text{for all}\ i\in I(x)\right\}.
\end{equation}
At a given point \((x,\lambda)\) and parameter $\mu$, the reduced gradient is defined as 
\begin{equation*}
	\proj{T(x)}{\nabla_x \Phi(x,\lambda,\mu)}.
\end{equation*}
Provided that the multipliers associated to the active bounds are of appropriate sign, the quantity $\|\proj{T(x)}{\nabla_x \Phi(x,\lambda,\mu)}\|$ is an appropriate criticality measure used in optimization algorithms for linearly constrained optimization such as MINOS~\cite{murtaghsaunders:1978} and LSNNO~\cite{tointtuyttens:1992}.
Projections on $T(x)$ can be computed in closed form. Indeed, assuming that $I(x) =  \left\{i_1,\ldots, i_p\right\}$ with $p < n-m$ and denoting by $Z \in \RR^{p\times n}$ the matrix whose row $k$ is the row $i_k$ of the $n\times n$ identity matrix, $T(x)$ is nothing but the null space of the block matrix 
\begin{equation}\label{eq:matrix_tangent_space}
	\tilde{A} = \begin{pmatrix}
	A \\Z
\end{pmatrix}.
\end{equation}
Let $\tilde{N}$ be a matrix whose columns form an orthonormal basis of the null space of $\tilde{A}$. By construction and our full rank Assumption~\ref{assumption:full_rank_A}, $\tilde{A}$ is also full rank. Then, the projection of a vector $v$ onto the null space of $\tilde{A}$ is given by $\tilde{P}v$ where
\[\tilde{P} = \tilde{N}\tilde{N}^\top,\]
There is no need to explicitly form the matrix $\tilde{N}$ because we can make use of the equivalent expression of the projection matrix $\tilde{P}$:
\begin{equation}\label{eq:minor_subpb_projector}
	\tilde{P} = I - \tilde{A}^\top\left(\tilde{A}\tilde{A}^\top\right)^{-1}\tilde{A}.
\end{equation}
Using~\eqref{eq:minor_subpb_projector}, one can compute the projection of a vector $v$, i.e. $\tilde{P}v$, by first solving, for an auxiliary variable $y$, the linear system
\begin{equation}\label{eq:proj_normal_eq_auxiliary_syst}
	\left(\tilde{A}\tilde{A}^\top\right)y = \tilde{A}v,
\end{equation}
and retrieve the projection by
\begin{equation*}
	v - \tilde{A}^\top y.
\end{equation*}
This approach, referred as the \textit{normal equations} approach~\cite{gould-etal:2001}, requires to solve the system~\eqref{eq:proj_normal_eq_auxiliary_syst}. This is done by using the Cholesky decomposition of $\tilde{A}\tilde{A}^\top$. The latter is well defined because the matrix $\tilde{A}\tilde{A}^\top$ is symmetric by construction and positive definite by Assumption~\ref{assumption:full_rank_A}. System~\eqref{eq:proj_normal_eq_auxiliary_syst} is reduced to two successive lower triangular systems. The factorization $\tilde{A}\tilde{A}^\top = \tilde{L}\tilde{L}^\top$, with $\tilde{L}$ lower triangular, does not need to be computed from scratch because the structure of $\tilde{A}$ implies a block structured Cholesky factor that involves the Cholesky decomposition of $AA^\top$. The latter is constant throughout the algorithm, so the factor $\tilde{L}$ can be computed in a relatively efficient manner, even though it potentially has to be updated at every new minor iteration. We expect it to be the case when these computations take place in the first outer iterations of Algorithm~\ref{algo:traulls}, but that as we progress toward a first-order critical point, the set of active bounds tends to stabilize and remain the same during the inner iterations. Details on this computation can be found in Appendix~\ref{appendix:chol_aug_matrix}.

This discussion justifies why our implementation uses condition
\begin{equation}\label{eq:criticality_cond_reduced_grad}
	\|\proj{T_k}{\nabla_x\Phi_k}\| \le \omega_k
\end{equation}
in Algorithm~\ref{algo:traulls} instead of~\eqref{eq:algo_traulls_criticality_test}.

\subsection{Inner minimization process}\label{sec:inner_iteration}

In this section, we describe the inner minimization process used to produce the outer iterates of Algorithm~\ref{algo:traulls}.
Since the current multipliers estimates $\lambda$ and the penalty parameter $\mu$ are fixed, the objective function for the inner minimization only depends on $x$ so we use the shorthand notation $\varphi : x\mapsto \Phi(x,\lambda,\mu)$.
Given a point $x_0 \in \Omega$, multipliers estimates $\lambda$, a penalty parameter $\mu$, we aim to find an approximate solution of the problem
\begin{equation}\label{eq:inner_itr_subpb}
	\begin{aligned}
		\min_x \quad & \varphi(x)\\
		\text{s.t.} \quad & x \in \Omega,
	\end{aligned}
\end{equation}
and such that
\begin{equation*}
	\|\proj{T(x)}{\nabla \varphi(x)}\| \le \omega,
\end{equation*}
for a tolerance $\omega > 0$. This task also involves an iterative process, for which the iterates will be indexed with subscript $k$ but are distinct from the outer iterates of Algorithm~\ref{algo:traulls}. At each iteration $k$, we compute a step $s_k$ and recursively update the iterate by $x_{k+1}=x_k+s_k$. The step is obtained after minimizing a model of the objective function $\varphi$ around $x_k$. We consider the quadratic model:
\begin{equation}\label{eq:al_quadratic_model}
	m_k(x) = \varphi_k + \inner{g_k}{x-x_k} + \dfrac{1}{2}\inner{x-x_k}{H_k(x-x_k)},
\end{equation}
where $\varphi_k:=\varphi(x_k)$, $g_k:=\nabla \varphi_k$ and $H_k$ is a symmetric approximation of the true Hessian $\nabla^2_{xx} \varphi_k$. Using the latter would impact negatively the performance of the algorithm because computing the second order terms~\eqref{eq:al_hessian_2nd_terms} requires too much time and storage in practice. We discuss the aspects of the methods relative to the choice of approximation in Section~\ref{subsec:hessian_approx}. 
Model~\eqref{eq:al_quadratic_model} is appropriate to describe the algorithm in terms of the iterate. When discussing subproblems, we prefer to emphasize the step and would then use the following model of the primal reduction $\varphi(x_k+s)-\varphi_k$, given by
\begin{equation*}
	q_k(s) = \dfrac{1}{2}\inner{s}{H_ks} + \inner{g_k}{s}.
\end{equation*}

One has $q_k(x-x_k)=m_k(x)-m_k(x_k)$.

We seek to compute the step $s_k$ as an approximate solution of the program

	\begin{align*}
		\min_{s} \quad& q_k(s)  \\
		\text{s.t.}  \quad & x_k+s \in \Omega \\ 
		& \|s\|_\infty \le \Delta_k,
	\end{align*}	

Vector $s$ denotes the unknown of the subproblem whose solution $s_k$ is the step and is used to compute the new iterate $x_{k+1}=x_k+s_k$.

Since $x_0$ is feasible, the constraints $x_k+s\in \Omega$ are satisfied at every iteration as long as the steps satisfy
\[As = 0, \qquad x_k-l \le s \le u-x_k,\]
for all $k$.

In our method, we also incorporate a trust-region strategy to control and assert the quality of a step. Therefore, we add to each subproblem the constraint $\|s\|_\infty \le \Delta_k$ with a radius $\Delta_k > 0$. This constraint reflects the domain on which we believe that the model well approximates the true function. Since the $\ell_\infty$ trust region constraint is equivalent to imposing $x_i \in [-\Delta_k,\Delta_k]$ for all $i$, we can rewrite this subproblem	
\begin{subequations}\label{pb:quadratic_subproblem} 
	\begin{align}
		\min_{s} \quad& q_k(s)  \\
		\text{s.t.}  \quad & As=0 \\
		& l_k \le s \le u_k,
	\end{align}	
\end{subequations}
where $l_k$ and $u_k$ are iteration-dependent bounds having components $(l_k)_i = \max\left(-\Delta_k,l_i-(x_k)_i\right)$ and $(u_k)_i = \min\left(\Delta_k,u_i-(x_k)_i\right)$ for all $i=1,\ldots,n$.

The success of an iteration and the update of the trust region are evaluated by the ratio
\begin{equation}\label{eq:tr_ratio}
	\rho_k = \dfrac{\varphi(x_k+s_k)-\varphi_k}{q_k(s_k)}.
\end{equation}
We follow the standard step acceptance criteria~\cite{conn-etal:2000}, which require constants $\eta_1, \eta_2, \gamma_1, \gamma_2$ such that 
\begin{equation}\label{eq:cond_constants_step_acceptance}
	0< \eta_1 \le \eta_2 < 1 \text{ and } 0< \gamma_1 \le \gamma_2 < 1.
\end{equation}
If $\rho_k  > \eta_1$, the step is accepted and the trust region is expanded. Otherwise, the step is rejected and the region reduced. A typical scheme to update the radius $\Delta_k$ would be to set
\begin{equation}\label{eq:tr_basic_update}
	\Delta_{k+1} \in \left\{\begin{aligned}
		& \left[\Delta_k,\infty\right) & &\text{if } \rho_k \ge \eta_2 \\
		& \left[\gamma_2\Delta_k, \Delta_k\right]  & &\text{if } \rho_k \in [\eta_1,\eta_2) \\
		& \left[\gamma_1\Delta_k, \gamma_2\Delta_k\right]  & &\text{if }  \rho_k<\eta_1
	\end{aligned}\right.
\end{equation}
The procedure employed to solve subproblem~\eqref{pb:quadratic_subproblem} is described in Algorithm~\ref{algo:trinner_iteration} and the latter is analyzed in the rest of this section. 
\begin{algorithm}
	\caption{Trust region inner iteration algorithm}
	\label{algo:trinner_iteration}
	\begin{algorithmic}[1]
		\Require initial point $x_0 \in \Omega$, radius $\Delta_0 >0$, constants $\eta_1, \eta_2, \gamma_1, \gamma_2$ satisfying conditions~\eqref{eq:cond_constants_step_acceptance}.
		\State set $k \gets 0$ 
		\Repeat{}
		\State compute the model $m_k$
		\State compute a step $s_k$ that sufficiently reduces the model
		\State compute the ratio $\rho_k$~\eqref{eq:tr_ratio}
		\If{$\rho_k > \eta_1$} set $x_{k+1}\gets x_k+s_k$ \Else{} set $x_{k+1} \gets x_k$ \EndIf
		\State set $\Delta_{k+1}$ according to~\eqref{eq:tr_basic_update}
		\State increment $k\gets k+1$
		\Until{$\|\proj{T_k}{g_k}\| \le \omega$}
	\end{algorithmic}
\end{algorithm}

\subsubsection{Computation of the step}\label{subsec:step_computation}

We consider the current feasible iterate $x_k$ and its associated approximate quadratic model $m_k$~\eqref{eq:al_quadratic_model}. The step is computed in two phases. We first compute a Cauchy step $s_k^C$ that ensures a decrease of the objective function sufficient to establish global convergence of the inner minimization algorithm. Next, we further minimize the objective function by exploring the subspace defined by the constraints active at the Cauchy point $x_k^C:=x_k+s_k^C$. This approach is of common use in gradient projection techniques~\cite[Chapter 16]{nocedalwright:2006} and there are different strategies to compute a Cauchy point~\cite{conn-etal:1988b,moretoraldo:1991,linmore:1999a}. Because of the structure of our constraints, we cannot directly project on the whole set $\Omega$ in practice but we can exploit prior knowledge on the active bounds. That is the reason why we compute our Cauchy step by finding the first local minimizer of the model along the projected gradient path
\begin{equation}\label{eq:projected_gradient_path}
	s_k(t)=\proj{\Omega}{x_k-tg_k}-x_k \text{ for } t\ge 0.
\end{equation}
Our procedure is an adaptation to the polyhedral case of algorithm SBMIN~\cite{conn-etal:1988b} used for the inner minimization phase of LANCELOT solver~\cite{conn-etal:1992}. 

Assume we have found a Cauchy step $s_k^C$. In order to have an efficient algorithm, we want the total step to achieve a better reduction than the Cauchy step. To do so, we build the next iterate $x_{k+1}$ after a finite sequence of $M$ minor iterates $x_{k,1},\ldots,x_{k,M+1}$. The sequence starts at the Cauchy point and ends at the next iterate, i.e. $x_{k}+s_k^C=x_{k,1}$ and $x_{k+1}=x_{k,M+1}$. This type of approach has shown to be effective for general optimization with bound constraints~\cite{linmore:1999a} and has also been incorporated into other AL algorithms for the inner loop minimization~\cite{conn-etal:1992,arreckx-etal:2016}.

Each minor iterate is defined after the previous one and is decomposed into 
\[x_{k,j+1}=x_{k,j}+w_{k,j},\]
where $w_{k,j}$ is a descent direction for the quadratic model $q_k$. For each minor iterate, we require
\begin{equation}\label{eq:minor_iterate_feasible}
	\begin{aligned}
		x_{k,j} \in \Omega, & & \|x_{k,j} - x_k \|_\infty \le \Delta_k, & & I(x_k^C) \subseteq I(x_{k,j}).
	\end{aligned}
\end{equation}
The first two conditions are merely that each minor iterate is feasible and the associated step lies within the trust region, while the last condition means that we can only add active bounds during this process. Note that in this context, a bound can become active with respect to the trust region and not only the original bounds on the variables.
We also require the sufficient decrease between two successive minor iterates
\begin{equation}\label{eq:minor_iterate_decrease}
	m_k(x_{k,j+1}) \le m_k(x_{k,j}), \qquad j = 1,\ldots,M.
\end{equation}
After each minor iteration, the corresponding step is $s_{k,j}:=x_{k,j}-x_k$.

The search direction $w_{k,j}$ is an approximate minimizer of the subproblem
\begin{subequations}\label{pb:minor_subpb}
	\begin{align}
		\min_w \quad & m_k(x_{k,j}+w) \\
		\text{s.t.} \quad & Aw=0 \label{subeq:minor_subpb_eq_cons} \\
		& w_i = 0, \qquad i \in I(x_{k,j}) \label{subeq:minor_subpb_fix_bounds}.
	\end{align}
\end{subequations}
At a minor iteration, the free variables, indexed by $\calF(x_{k,j})$, are implicitly subject to the bounds 
\begin{equation}\label{subeq:minor_subpb_bounds}
	\left(l_{(k,j)}\right)_i \le w_i \le \left(u_{(k,j)}\right)_i, \qquad i \in \calF(x_{k,j})
\end{equation}
where $l^{(k,j)}:= l_k - s_{k,j}$ and $u^{(k,j)}:= u_k - s_{k,j}$.

The minor subproblem~\eqref{pb:minor_subpb} is solved by applying the projected conjugate gradient~\cite{gould-etal:2001} method with the previous minor iterate $x_{k,j}$ as a starting point. Three termination cases can occur. First, a direction can be generated such that one of the bounds~\eqref{subeq:minor_subpb_bounds} is violated for one of the free variables. When this happens, the search direction is scaled so that the associated component lies at the bound and the CG iterations stop. The second case occurs when we generate a direction of negative curvature and is handled similarly to the first one, i.e. we modify the direction so that all the variables remain within their bounds. The third case is the normal termination one and occurs when a local minimizer -- with respect to a given tolerance -- has been found. In all cases, we return the obtained search direction $w_{k,j}$, perform the projected line search to compute the point $x_{k,j+1}$ such that~\cref{eq:minor_iterate_feasible,eq:minor_iterate_decrease} are verified. The handling of the termination cases differs in the continuation of the procedure. The last two cases (negative curvature and local minimality), cause the stopping of the minor iterates loop. On the contrary, if the CG iterations stopped because a bound was hit, we go back to solving~\eqref{pb:minor_subpb} using projected conjugate gradient method with $x_{k,j+1}$ as a starting point and $I(x_{k,j+1}) \supset I(x_{k,j})$ as a new set of fixed components.

Each iteration of the conjugate gradient method applied to~\eqref{pb:minor_subpb} requires computing projections onto the null space of the constraints~\eqref{subeq:minor_subpb_eq_cons}--\eqref{subeq:minor_subpb_fix_bounds}. We denote the associated projection operator, in this case a $n\times n$ matrix, by $\tilde{P}$. The specification of the projected conjugate gradient method applied to solving~\eqref{pb:minor_subpb} is outlined in Algorithm~\ref{algo:projected_cg_method}.
\begin{algorithm}
	\caption{The projected conjugate gradient method applied to~\eqref{pb:minor_subpb}}\label{algo:projected_cg_method}
	\begin{algorithmic}[1]
		\State\label{line:pcg_initial_projection}Set $w \gets 0,\ r \gets H_k(x_{k,j}-x_k)+g_k,\ v\gets \tilde{P}r,\ p \gets -v$
		\State Set tolerance $\varepsilon_{cg} \gets \kappa_{cg}\|v\|$
		\Repeat
		\If{$\inner{p}{H_kp}\le 0$}
		\State Set $w^+ \gets w + \gamma p$ where $\gamma$ is the smallest factor such that $w^+_i \in \{l^{(k,j)}_i,u^{(k,j)}_i\}$ 
		\State for $i \in \calF(x_{k,j})$
		\State \textbf{STOP}
		\EndIf
		\State Set $\alpha \gets \inner{r}{v} / \inner{p}{H_kp}$
		\State Set $w^+ \gets w+\alpha p$
		\If{$w^+$ violates a bound}
		\State $w^+ \gets w + \gamma p$ where $\gamma$ is the smallest factor such that $w^+_i \in \{l^{(k,j)}_i,u^{(k,j)}_i\}$
		\State for $i \in \calF(x_{k,j})$
		\State \textbf{STOP}
		\EndIf
		\State Set $r^+ \gets r + \alpha H_kp$
		\State\label{line:pcg_iteration_projection} Set $v^+\gets \tilde{P}r^+$
		\If{$\sqrt{\inner{r^+}{v^+}} < \varepsilon_{cg}$} \textbf{STOP}
		\EndIf
		\State Set $\beta \gets \inner{r^+}{v^+} / \inner{r}{v}$ 
		\State Set $p \gets -v^+ + \beta p$
		\State Set $w\gets w^+,\ r \gets r^+,\ v \gets v^+$
		\Until{$2(n-m-|I(x_{k,j})|)$ iterations have been done} 
		\State{}
		\Return $w^+$
	\end{algorithmic}
\end{algorithm} 
The projections occurring at lines~\ref{line:pcg_initial_projection}~and~\ref{line:pcg_iteration_projection}, and the associated operators, are computed with the normal equations approach described in previous section.

An interesting feature would be to complement each search direction by a steplength computed by a projected search~\cite{moretoraldo:1991,linmore:1999a}. This allows to add more than one constraint at each minor iteration to the active set. The downside is that it would also require to compute the projection of the gradient direction~\eqref{eq:projected_gradient_path} on the feasible set $\Omega$ at every new trial value of the steplength. When $\Omega$ only contains bound constraints, the projection is trivial and cheap to compute. The linear equality case is more costly but can still be handled efficiently by a \textit{normal equations} or \textit{augmented system} approach~\cite{gould-etal:2001}. When $\Omega$ is of the form~\eqref{eq:linear_constraints}, the projection is less trivial and requires to solve the associated minimum-distance quadratic program with a dedicated solver for quadratic programming, which we prefer to avoid.

We stop the minor-iterations procedure when the reduced gradient associated to the current active set is small enough. More formally, assume we performed the $j+1$ minor minimizations and let $T_{k,j}$ be the tangent space at $x_{k,j}$ with respect to the active bounds in $I(x_{k,j})$. We stop the minor iteration loop whenever the following inequality is satisfied:
\begin{equation*}
	\left\Vert \proj{T_{k,j}}{\nabla m_k(x_{k,j+1})} \right\Vert \le \kappa_{mlt} \left\Vert \proj{T_{k,j}}{\nabla m_k(x_k)}\right\Vert,
\end{equation*}
with $\kappa_{mlt} \in (0,1)$. This inequality estimates if there is relative progress that can be made in the tangent subspace spanned by the free variables. 

\subsubsection{Hessian approximation}\label{subsec:hessian_approx}

We now discuss how the Hessian of the AL is iteratively updated when we define the quadratic model of iteration $k$ in Algorithm~\ref{algo:trinner_iteration}. To set up the context and notations of this paragraph, we wish to approximate the true Hessian
\[\nabla_{xx}^2 \varphi(x_k) = J_k^\top J_k + \mu C_k^\top C_k + S_k,\]  
by a symmetric matrix $H_k$.
We remind that the need for an approximation mainly comes from the fact that evaluating the second order terms in $S_k$ requires too much time and storage to be done in practice, especially for problems with a large number of variables and residuals.

The first approximation we can think of, and the simplest one, is obtained after merely linearizing the residuals and constraints in expression~\eqref{eq:al}, which gives
\begin{equation}\label{eq:hessian_gn_approx}
	H_k = J_k^\top J_k + \mu C_k^\top C_k.
\end{equation}
We also call it the Gauss-Newton (GN) approximation, in reference to its counterpart in the unconstrained case. On the one hand, this approximation is very convenient and cheap as it involves already available first-order derivatives, since they are required to evaluate the gradient. Also, when the Jacobians are full rank, the resulting matrix is positive-definite which guarantees the convexity of subproblems of the form~\eqref{eq:inner_itr_subpb}. On the other hand, it is known to have a slower rate of convergence on problems with non zero residuals at the solution. This downside is amplified when we are to approximate the Hessian of the Lagrangian, or the AL in our case. Indeed, setting $S_k$ to the zero matrix not only neglects the contribution of the residuals to the curvature of the model, but also neglects the contribution of the Lagrange multipliers, for which there are no reasons to equal zero.

We thus need to take into account the full Hessian to form an approximation. Looking at the literature on this subject, one can observe that there is a variety of approaches focusing on the unconstrained case, as it is a major challenge in improving the efficiency of algorithms for problems with non zero residuals at the solution. Nevertheless, relevant parallels can be drawn with the AL situation, or the constrained case in general, since these studies explore techniques to deal with second order terms. Because the first-order terms are readily available and provide curvature information, only the second-order components need to be approximated. Therefore, we look for an approximation of the form 
\begin{equation}\label{eq:hessian_full_approx}
	H_k = B^{GN}_k + B_k.
\end{equation}
where $B^{GN}_k$ is the right hand side of~\eqref{eq:hessian_gn_approx} and $B_k$ is a symmetric approximation of $S_k$. 

Structured approximations of the AL Hessian have been studied for general objective functions~\cite{tapia:1988} but our approach is closer to the one employed in~\cite{li-etal:2002}. In the latter, the authors base their approximation on a secant equation derived from a heuristic initially introduced in the unconstrained case~\cite{biggs:1977} that they adapted to approximate the Hessian of the Lagrangian and use it in a SQP algorithm for constrained nonlinear least-squares. This strategy is also at the heart of the SQN method from the popular package for unconstrained nonlinear least-squares NL2SOL~\cite{dennisetal:1981,dennis-etal:1989}. The idea is the following. If we were to approximate each matrix term of the sum~\eqref{eq:al_hessian_2nd_terms}, we would have
\begin{equation}\label{eq:sum_hessian_approximations}
	B_{k+1} = \sum_{i=1}^{n_r}  r_i(x_{k+1}) B^{r_i}_{k+1} + \sum_{i=1}^{n_c} \bar{\lambda}_i(x_{k+1},\lambda,\mu) B^{c_i}_{k+1},
\end{equation}
with $B^{r_i}_{k+1} \approx \nabla^2r_i(x_{k+1})$ and $B^{c_i}_{k+1} \approx  \nabla^2 c_i(x_{k+1})$. To get an accurate approximation, given a step $s_k$, it is reasonable to require
\begin{equation}\label{eq:components_secant_equations}
	B^{r_i}_{k+1}s_k = \nabla r_i(x_{k+1}) - \nabla r_i(x_k) \qquad \text{and} \qquad B^{c_i}_{k+1}s_k = \nabla c_i(x_{k+1}) - \nabla c_i(x_k),
\end{equation}
i.e. that each Hessian maps the change in the variables to the change in the gradients. Noticing that, for each $i$, $\nabla r_i(x_{k+1}) - \nabla r_i(x_k)$, resp. $\nabla c_i(x_{k+1}) - \nabla c_i(x_k)$, is the $i$-th row of $\left(J_{k+1}-J_k\right)^\top$, resp. $\left(C_{k+1}-C_k\right)^\top$, summing over the residuals and constraints indices all the terms of~\eqref{eq:components_secant_equations} gives, by~\eqref{eq:sum_hessian_approximations}, the structured secant equation
\begin{equation}\label{eq:structured_secant_eq}
	B_{k+1}s_k = \left(J_{k+1}-J_k\right)^\top r_{k+1} + \left(C_{k+1}-C_k\right)^\top\bar{\lambda}_{k+1}.
\end{equation}
To follow the conventional notations of quasi-Newton methods, we will denote the right hand side of~\eqref{eq:structured_secant_eq} by $\tilde{y}_k$ to insist on its distinction with the standard term $y_k:= g_{k+1}-g_k$. Note that, as in SQN methods for unconstrained least-squares~\cite{lucksan-etal:2019}, one has
\[B^{GN}_ks_k + \tilde{y}_k = y_k,\]
so the resulting approximation satisfies the standard secant equation $H_ks_k = y_k$.

We can now derive an update formula from equation~\eqref{eq:structured_secant_eq}. In our implementation, we apply the SR1 update
\begin{equation}\label{eq:structured_sr1_update}
	B_{k+1} = B_k + \dfrac{(\tilde{y}_k-B_ks_k)(\tilde{y}_k-B_ks_k)^\top}{\inner{\tilde{y}_k-B_ks_k}{s_k}},
\end{equation}
if $\inner{\tilde{y}_k-B_ks_k}{s_k} \neq 0$. When the denominator in~\eqref{eq:structured_sr1_update} is zero, we simply set $B_{k+1}=B_k$. To avoid numerical errors, we apply the update when
\begin{equation}\label{eq:structured_sr1_safeguard}
	\left|\inner{\tilde{y}_k-B_ks_k}{s_k}\right| \ge \kappa_{sds}\|s_k\|\|\tilde{y}_k-B_ks_k\|,
\end{equation}
for $\kappa_{sds} \in (0,1)$. Inequality~\eqref{eq:structured_sr1_safeguard} is one of the most common safeguards used in SR1 methods. 

The choice of the SR1 method is motivated by its robustness in the least-squares context, as the exhaustive benchmarks in~\cite{lucksan-etal:2019} show, and its tendency to better approximate the true Hessian with each iteration~\cite{conn-etal:1991a}. The approximation could thus be indefinite, which could be considered as a weakness compared to BFGS or DFP updates, that are guaranteed to be positive definite as long as it is the case for the initial approximation $B_0$. However, as it is highlighted in Algorithm~\ref{algo:projected_cg_method}, the use of the conjugate gradients method in a trust region framework handles, and even exploits, the potential non-convexity of the subproblems. Moreover, this update does not require a curvature condition $\inner{s_k}{y_k} > 0$ to be satisfied. In other works on the constrained case, SQN methods are used to approximate the projected Hessian of the Lagrangian \cite{tjoabiegler:1991}, or of a merit function~\cite{amiribartels:1989}, in the null space of the active constraints. The BFGS update is then preferred because this matrix is, under standard assumptions, positive semidefinite, according to the second-order necessary conditions.

We end our description of the Hessian approximation by discussing hybrid updates, a feature commonly used in SQN methods for nonlinear least-squares algorithms, either in the unconstrained~\cite{dennisetal:1981,albaalifletcher:1985,fletcherxu:1987} or constrained~\cite{tjoabiegler:1991,li-etal:2002} case. The idea is to test at every iteration if the problem has small or large residuals at the solution. In the latter case, a structured update is used to increase the accuracy of the approximated Hessian whereas in the former, the GN approximation is employed, as it is more likely to correspond to the true Hessian. Note that the update, such as~\eqref{eq:structured_sr1_update}, is still computed at every iteration to continue accumulating curvature information. We use a strategy inspired from~\cite{fletcherxu:1987} adapted to the constrained case in~\cite{tjoabiegler:1991}, that computes the relative reduction
\begin{equation*}
	\zeta_k = \dfrac{\varphi(x_{k})-\varphi(x_{k+1})}{\varphi(x_{k})},
\end{equation*}
and updates the approximated Hessian based on the rule
\begin{equation}\label{eq:hybrid_update_strategy}
	H_{k+1} = \left\{\begin{aligned}
		&B^{GN}_{k+1} & &\text{if } \zeta_k \le \kappa_{hyb} \\
		&B^{GN}_{k+1} + B_{k+1} & & \text{otherwise}, 
	\end{aligned}\right.
\end{equation}
where $\kappa_{hyb}$ is a parameter in $(0,1)$. This update is used with $\kappa_{hyb}=0.1$ in~\cite{tjoabiegler:1991}, where the AL plays the role of a merit function. As we will see in Section~\ref{sec:numerical_experiments}, it has an important impact on the performance of the algorithm.

\subsubsection{Summary and implementation details}\label{subsec:implementation_details}

We end our discussion by summarizing the relations between outer, inner and minor iterates and the default values for the different constants exposed in this section.
\begin{itemize}
	\item The outer iterates $x_K$ are the main iterates of the algorithm and are approximate minimizers of the AL
	\item Each outer iterate is formed after a sequence of inner iterates $(x_k)_k$, linked by $x_{k+1}=x_k+s_k$ where $s_k$ is an approximate solution of the QP~\eqref{pb:quadratic_subproblem}
	\item Each inner iterate is formed after a sequence of $M$ minor iterates $x_{k,j}$ linked by $x_{k,j+1}=x_{k,j}+w_{k,j}$ where $w_{k,j}$ is a descent direction obtained by applying the projected conjugate gradient Algorithm~\ref{algo:projected_cg_method}
\end{itemize} 
The relations between the three levels of iterations are illustrated in Figure~\ref{fig:iteration_hierarchy}.
\begin{figure}[ht]
	\centering
	\begin{tikzpicture}[
		>=Stealth,
		every node/.style={font=\small},
		outerbox/.style={draw=blue!60, fill=blue!5, rounded corners=8pt, thick, inner sep=10pt},
		innerbox/.style={draw=teal!70, fill=teal!5, rounded corners=6pt, thick, inner sep=8pt},
		minorbox/.style={draw=orange!70, fill=orange!5, rounded corners=5pt, thick, inner sep=6pt},
		stepbox/.style={draw=gray!60, fill=gray!8, rounded corners=3pt, inner sep=5pt, align=center, font=\footnotesize},
		]
		% Outer iteration frame
		\node[outerbox, minimum width=13.5cm, minimum height=10.5cm, label={[font=\small\bfseries, anchor=north]north:Outer iterations (Algorithm~\ref{algo:traulls})}] (outer) {};
		\node[font=\footnotesize, anchor=north] at ([yshift=-4mm]outer.north) {Approx.\ minimize AL $\Phi$, update $\lambda$, $\mu$, tolerances};
		
		% Inner iteration frame
		\node[innerbox, minimum width=12cm, minimum height=8.2cm, label={[font=\small\bfseries, anchor=north]north:Inner iterations (Algorithm~\ref{algo:trinner_iteration})}] (inner) at ([yshift=-6mm]outer.center) {};
		\node[font=\footnotesize, anchor=north] at ([yshift=-4mm]inner.north) {Trust region steps: $x_{k+1}=x_k+s_k$};
		
		% Left column: Cauchy + TR update + Hessian update + convergence check
		\node[stepbox, minimum width=3.8cm] (cauchy) at ([xshift=-2.2cm, yshift=1.8cm]inner.center) {Cauchy step\\[1pt]{\scriptsize Projected gradient path}};
		\node[stepbox, minimum width=3.8cm, below=8mm of cauchy] (trupdate) {Trust region update\\[1pt]{\scriptsize Accept/reject via $\rho_k$}};
		\node[stepbox, minimum width=3.8cm, below=8mm of trupdate] (hessupdate) {Hessian update\\[1pt]{\scriptsize Structured SR1}};
		\node[stepbox, minimum width=3.8cm, below=8mm of hessupdate] (convcheck) {$\|P_{T_k}[\nabla\varphi_k]\| \le \omega$\,?};
		
		% Minor iteration frame
		\node[minorbox, minimum width=4.5cm, minimum height=6.2cm, right=8mm of cauchy, anchor=north west, yshift=8mm, label={[font=\small\bfseries, anchor=north]north:Minor iterations}] (minor) {};
		\node[font=\footnotesize, anchor=north] at ([yshift=-4mm]minor.north) {$x_{k,j+1}=x_{k,j}+w_{k,j}$};
		
		% Inside minor: PCG + termination
		\node[stepbox, minimum width=3.5cm] (pcg) at ([yshift=6mm]minor.center) {Projected CG\\[1pt]{\scriptsize Algorithm~\ref{algo:projected_cg_method}}};
		\node[stepbox, minimum width=3.5cm, below=7mm of pcg] (boundhit) {Bound hit};
		\node[stepbox, minimum width=3.5cm, below=7mm of boundhit] (negcurv) {Neg.\ curvature / converged};
		
		% Arrows left column
		\draw[->, thick, gray] (cauchy) -- (trupdate);
		\draw[->, thick, gray] (trupdate) -- (hessupdate);
		\draw[->, thick, gray] (hessupdate) -- (convcheck);
		
		% Arrow Cauchy -> minor
		\draw[->, thick, gray] (cauchy.east) -- ([xshift=-2mm]pcg.west|-cauchy.east) -- ([xshift=-2mm]pcg.west) -- (pcg.west);
		
		% Arrows inside minor
		\draw[->, thick, gray] (pcg) -- node[right, font=\scriptsize]{or} (boundhit);
		\draw[->, thick, gray] (boundhit) -- node[right, font=\scriptsize]{or} (negcurv);
		
		% Restart loop for bound hit
		\draw[->, thick, gray, dashed] (boundhit.east) -- ++(6mm,0) |- node[right, font=\scriptsize, pos=0.25]{restart} (pcg.east);
		
	\end{tikzpicture}
	\caption{Hierarchy of the three iteration levels in the AL algorithm. The outer loop updates the Lagrange multipliers and penalty parameter. Each outer iteration triggers a sequence of trust-region inner iterations, and each inner iteration computes its step via a Cauchy point followed by minor iterations using the projected conjugate gradient method.}
	\label{fig:iteration_hierarchy}
\end{figure}
Our intention with the work presented in this paper was to conceive an algorithm able to handle directly linear constraints of the form~\eqref{eq:linear_constraints} but we also accept problems where linear constraints are only bounds and then
\[\Omega = \left\{x \in \RR^n \ |\ l \le x \le u\right\}.\]
The framework we have presented remains valid for this formulation. The only practical difference is in the computation of projections. In the bound constrained  case, the projection of a vector $x$ on the set $\Omega$ has components
\begin{equation}\label{eq:projection_bounds}
	\left(\proj{\Omega}{x}\right)_i = 
	\left\{ \begin{aligned}
		& l_i & &\text{if } x_i < l_i \\
		& x_i & &\text{if } x_i \in [l_i,u_i] \\
		& u_i & &\text{if } x_i > u_i
	\end{aligned} \right.
\end{equation}
In our implementation, the only practical difference with the polyhedral case is that we can use directly the norm of the projected gradient
\[\left\| x_K - \proj{\Omega}{x_K-\nabla_x\Phi_K}\right\|,\]
as a criticality measure.

The default values for tolerances and constants of Algorithm~\ref{algo:traulls} are
\[ \mu_0 = 10,\ \kappa_\omega = \beta_\omega = \eta = \omega = 1,\ \kappa_\eta = 0.1,\ \beta_\eta = 0.9,\ \tau = 100,\  \omega_*=\eta_* = 10^{-7}.\]
The constants relative to the inner minimization process of Algorithm~\ref{algo:trinner_iteration} are set to the subproblem criticality tolerance, i.e. at outer iteration $K$, we set
\[ \kappa_{cg} = \kappa_{mlt} = \omega_K.\]
The value of constant $\kappa_{sds}$ used safeguard~\eqref{eq:structured_sr1_safeguard} is set to the square root of the relative double precision.

At the start of each outer iteration, we set our initial approximation of the second order terms of the Hessian to $B_0 = 0$ and the constant used to monitor the hybrid switching scheme~\eqref{eq:hybrid_update_strategy} of the Hessian is set to $\kappa_{hyb}=0.1$.

We now discuss the trust region handling. The initial radius value for Algorithm~\ref{algo:trinner_iteration} is set to
\begin{equation*}
	\Delta_0 = 0.1\| g_0\|_\infty.
\end{equation*}
The mechanism to update the trust region given in~\eqref{eq:tr_basic_update} still leaves important flexibility. We choose to follow a standard strategy similar to the one exposed in~\cite[][Chapter 17]{conn-etal:2000}. We also added a refinement to handle the case where the ratio $\rho_k$ is negative, which can happen when there is very poor agreement between the function and the model. In such cases, we reduce more severely the radius by taking
\[\Delta_{k+1} = \gamma_1 \Delta_k.\]
The complete update rule for the trust region radius is then
\begin{equation*}
	\Delta_{k+1} = \left\{\begin{aligned}
		& \max\left(\alpha_2 \|s_k\|_\infty,\Delta_k\right) & &\text{if } \rho_k \ge \eta_2 \\
		& \Delta_k &  &\text{if } \rho_k \in [\eta_1,\eta_2) \\
		& \alpha_1\|s_k\|_\infty & &\text{if } \rho_k \in [0,\eta_1) \\
		& \min\left(\alpha_1 \|s_k\|_\infty,\gamma_1\Delta_k\right) & &\text{if } \rho_k < 0,
	\end{aligned}\right.
\end{equation*}
with constant values
\[\alpha_1 = 0.25,\ \alpha_2 = 2.5,\ \eta_1 = 0.25,\ \eta_2 = 0.75,\ \gamma_1 = 0.0625.\]
